We consider learning from demonstrations in deterministic planning problems.
These problems are goal-based and object-centric, with continuous states and hybrid discrete-continuous actions.
Formally, an \emph{environment} is a tuple $\langle \types, d, \C, f, \Psi_G \rangle$, and is associated with a distribution $\T$ over \emph{tasks}, where each task $T \in \T$ is a tuple $\langle \O, x_0, g \rangle$.

$\types$ is a finite set of object \emph{types}, and the map $d: \types \to \mathbb{N}$ defines the dimensionality of the real-valued feature vector for each type. Within a task, $\O$ is an \emph{object set}, where each object has a type drawn from $\types$; this $\O$ can (and typically will) vary between tasks.
$\O$ induces a state space $\X_\O$ (going forward, we simply write $\X$ when clear from context). A \emph{state} $x \in \X$ in a task is a mapping from each $o \in \O$ to a feature vector in $\mathbb{R}^{d(\text{type}(o))}$; $x_0$ is the initial state of the task.

$\C$ is a finite set of \emph{controllers}. A controller $C((\type_1, \dots, \type_v), \Theta) \in \C$ can have both discrete typed parameters $(\type_1, \ldots, \type_v)$ and a continuous real-valued vector of parameters $\Theta$. For instance, a controller \texttt{Pick} for picking up a block might have one discrete parameter of type \texttt{block} and a $\Theta$ that is a placeholder for a specific grasp pose. The controller set $\C$ and object set $\O$ induce an action space $\A_\O$ (going forward, we write $\A$ when clear). An \emph{action} $a \in \A$ in a task is a controller $C \in \C$ with both discrete and continuous arguments: $a = C((o_1, \ldots o_v), \theta)$, where the objects $(o_1, \ldots o_v)$ are drawn from the object set $\O$ and must have types matching the controller's discrete parameters $(\type_1, \ldots, \type_v)$. 
Transitions through states and actions are governed by $f: \X \times \A \to \X$, a known, deterministic transition model that is shared across tasks.

A \emph{predicate} $\psi$ is characterized by an ordered list of types $(\type_1, \dots, \type_m)$ and a lifted binary state classifier $c_\psi : \X \times \O^m \to \{\text{true}, \text{false}\}$, where $c_\psi(x, (o_1, \dots, o_m))$ is defined only when each object $o_i$ has type $\type_i$.
For instance, the predicate \mbox{\texttt{Holding}} may, given a state and two objects, robot and block, describe whether the block is held by the robot in this state.
A \emph{lifted atom} is a predicate with typed variables (e.g., \texttt{Holding(?robot, ?block)}).
A \emph{ground atom} $\ground{\psi}$ consists of a predicate $\psi$ and objects $(o_1, \dots, o_m)$, again with all $\text{type}(o_i) = \type_i$ (e.g., \texttt{Holding(robby, block7)}).
Note that a ground atom induces a binary state classifier $c_{\ground{\psi}} : \X \to \{\text{true}, \text{false}\}$, where $c_{\ground{\psi}}(x) \triangleq c_{\psi}(x, (o_1, \dots, o_m))$.

$\Psi_G$ is a small set of \emph{goal predicates} that we assume are given and sufficient for representing task goals, but insufficient practically as standalone state abstractions. Specifically, the goal $g$ of a task is a set of ground atoms over predicates in $\Psi_G$ and objects in $\O$.
A goal $g$ is said to \emph{hold} in a state $x$ if for all ground atoms $\ground{\psi} \in g$, the classifier $c_{\ground{\psi}}(x)$ returns true.
A solution to a task is a \emph{plan} $\pi = (a_1, \ldots, a_n)$, a sequence of actions $a \in \A$ such that successive application of the transition model $x_i = f(x_{i-1}, a_i)$ on each $a_i \in \pi$, starting from $x_0$, results in a final state $x_n$ where $g$ holds.

The agent is provided with a set of \emph{training tasks} from $\T$ and a set of demonstrations $\D$, with \emph{one demonstration per task}.
We assume action costs are unitary and demonstrations are near-optimal.
Each demonstration consists of a training task $\langle \O, x_0, g \rangle$ and a plan $\pi^*$ that solves the task.
Note that for each $\pi^*$, we can recover the associated state sequence starting at $x_0$, since $f$ is known and deterministic.
The agent's objective is to \emph{efficiently} solve held-out tasks from $\T$ using anything it chooses to learn from $\D$.
% In other words, the agent should produce \emph{good} solutions \emph{quickly}.
% In our experiments, to assess generalization, we evaluate the agent on tasks that involve more objects, require longer action sequences, and have larger goal expressions than the training tasks.